\begin{Exercice}[][15 min]
On considère la fonction numérique $f$ définie sur $\mathbb{R}$ par : \quad
$f\big(x\big) = \dfrac{2x}{x^2 + 1}$
\vspace{0.2cm} \\
On note $\big(\mathcal{C}_f\big)$ sa courbe représentative dans un repère orthonormé $\big(O; \vect{i}; \vect{j}\big)$.

\begin{enumerate}
    \item Déterminer $D_f$, puis étudier la parité de $f$.
    \item Calculer $f\big(0\big)$, $f\big(1\big)$ et $f\big(2\big)$.

    \item Soient $a$ et $b$ deux réels distincts de l'intervalle $[0; +\infty[$.
    \begin{enumerate}
        \item Montrer que :
        $T\big(a; b\big) = \dfrac{2\big(1 - ab\big)}{\big(a^2 + 1\big)\big(b^2 + 1\big)}$
        \item Étudier la monotonie de $f$ sur $[0; 1]$ et sur $[1; +\infty[$.
        \item En déduire la monotonie de $f$ sur $[-1; 0]$ et sur $]-\infty; -1]$.
        \item Dresser le tableau de variations de $f$ sur $\R$.
    \end{enumerate}
    

    \item 
    \begin{enumerate}
        \item Montrer que pour tout réel $x \in \mathbb{R}$, on a : $1 - f\big(x\big) = \dfrac{\big(x - 1\big)^2}{x^2 + 1}$
        \item En déduire que pour tout $x \in \mathbb{R}$, on a : $f\big(x\big) \leqslant 1$
    \end{enumerate}
    \item Résoudre dans $\mathbb{R}$ l'équation $f\big(x\big) = 1$
          En déduire que le nombre $1$ est le maximum de $f$ sur $\mathbb{R}$.
    \item En utilisant la parité de $f$, montrer que $-1$ est le minimum de $f$ sur $\mathbb{R}$ et préciser en quel point il est atteint.
\end{enumerate}
\end{Exercice}
